% ============================================================
%  01_introduccion.tex — Capítulo 1: Introducción
% ============================================================

\chapter{Introducción}
\label{cap:introduccion}

% ------------------------------------------------------------------
% TIP: este es el esqueleto. Rellena cada sección con tu contenido.
% ------------------------------------------------------------------

\section{Antecedentes}
\label{sec:antecedentes}

\textit{[Contexto general del problema. ¿Qué se sabe del tema? ¿Qué avances existen
en agricultura de precisión, invernaderos, sensórica e Internet de las cosas?]}

\section{Planteamiento del problema}
\label{sec:planteamiento}

\textit{[Describe el problema concreto que la tesis pretende resolver. ¿Qué dificultad
existe al gestionar un invernadero? ¿Por qué es necesario un enfoque basado en
ontologías / sistemas inteligentes?]}

\section{Justificación}
\label{sec:justificacion}

\textit{[¿Por qué es importante este trabajo? Impacto académico, económico y social.]}

\section{Objetivos}
\label{sec:objetivos}

\subsection{Objetivo general}
\label{subsec:obj-general}

\textit{[Enunciar el objetivo general en una frase.]}

\subsection{Objetivos específicos}
\label{subsec:obj-especificos}

\begin{enumerate}
    \item \textit{[Objetivo específico 1. Ej.: diseñar una ontología del dominio del invernadero.]}
    \item \textit{[Objetivo específico 2.]}
    \item \textit{[Objetivo específico 3.]}
\end{enumerate}

\section{Alcance y limitaciones}
\label{sec:alcance}

\textit{[¿Qué cubre y qué no cubre el trabajo?]}

\section{Estructura de la tesis}
\label{sec:estructura}

\textit{[Breve descripción de cómo está organizado el documento (capítulo por capítulo).]}
